% Target: 6--8 pages once exposition is added. See ../PROJECT_REQUIREMENTS.md.
% Compact metrics, matrices and ROC/PR restored here; expanded versions in supplement.tex.
\documentclass[sigconf,nonacm]{acmart}
\setcopyright{none}
\settopmatter{printacmref=false,printccs=false,printfolios=true}
\acmDOI{}\acmISBN{}
\usepackage{booktabs}
\usepackage{graphicx}
\usepackage{adjustbox}
\newcommand{\draftspace}{\mbox{}\par\smallskip}
\newcommand{\protocolplot}[2]{%
\begin{figure*}[p]
\centering
\includegraphics[width=\textwidth,height=.43\textheight,keepaspectratio]{figures/#1-curves.pdf}
\caption{#2. ROC (top), precision--recall (bottom); mismatch-positive, scored pairs only.}
\Description{Per-protocol curves with model legends and separate evaluation panels.}
\label{fig:protocol-#1}
\end{figure*}}
\title{Checking Semantic Correspondence Between Natural-Language and Lean Proofs}
\author{Taylor Belcher}
% TODO: affiliation/email if required.
\renewcommand{\shortauthors}{Belcher}
\begin{document}
\begin{abstract}
% TODO: Taylor's abstract.
\end{abstract}
\keywords{Lean, proof correspondence, informalization, LoRA}
\maketitle

\section{Introduction}
% Budget: about 0.6 page. Restate project-proposal-tbelcher.tex briefly.
\subsection{Lean Validity, Theorem Equivalence, and Argument Fidelity}
\draftspace
% TODO: motivate the checker using the Navier--Stokes critique and Xena blog;
% attribute reported discrepancies rather than asserting independently verified ones.
\subsection{Research Question and Contributions}
\draftspace
% TODO: dataset/evaluation challenges; informalization as checker component and
% independently useful aid to human reading. No new formalizer is the objective.

\section{Related Work}
% Budget: about 0.4 page. Cite collected papers/blogs selectively.
\subsection{Faithful Formalization and Reported Translation Failures}
\draftspace
% TODO: Pistis, FaithformBench, Navier--Stokes critique, Xena, George blog.
\subsection{Informalization and Intermediate Representations}
\draftspace
% TODO: English, controlled English, Verbose Lean, Natural Lean; retain their
% potential without claiming that our code intermediaries were verified successes.

\section{Preliminary Experiments: The Dataset Challenge}
% Brief historical exposition; compact numerical evidence is restored below.
\subsection{Statement Judges, Pseudo-Labels, and Proof-Swap Training}
\draftspace
% TODO: 8B statement-only first/second pass; 14B random/related proof swaps.
% Distinguish source-assumed positives, synthetic negatives and pseudo-labels.
\subsection{Random Forests, Same-Theorem Diagnostics, and Intermediaries}
\draftspace
% TODO: lexical shortcuts, tiny constructed diagnostic, missing pair labels.
% Also acknowledge calibration and generation limits: not every limitation
% established a dataset-caused failure. Different protocols are not comparable.
\subsection{Lessons for the Final Experiment}
\draftspace
% TODO: why separate translation supervision from correspondence judgments.
\input{tables/main-metrics}
\input{tables/main-confusion}
\protocolplot{statements}{8B statement correspondence}
\protocolplot{random-swap}{14B random-swap training and same-theorem diagnostic}
\protocolplot{forest}{Paired English--Lean random forest}
\protocolplot{related-swap}{14B related-swap training and evaluation}
\protocolplot{intermediaries}{Original-14B translation intermediaries; unequal coverage}
\protocolplot{pilot32}{Original 32B provisional mixed-source pilot}
\protocolplot{english-control}{English-only forest: secondary shortcut control}
\protocolplot{masked-control}{Lean-proof-withheld secondary control}

\section{Informalization Pipeline: Data and Methods}
% Budget: about 1.5 pages; together with results, the majority of main-body space.
% Preserve rubric coverage: Data 10%, Methods 25%.
\subsection{Herald Translation Pairs and Grouped Held-Out Splits}
\draftspace
% TODO: original Lean/English pairs in reverse; 4096/64/64 selected examples;
% full proof extraction, token exclusions, duplicate/family separation and limits.
\subsection{Stage One: Training Qwen3-14B to Explain Lean Proofs}
\draftspace
% TODO: fresh BF16 base + LoRA; English completion-only loss; Lean statement,
% body and context as input, no original English available to the translator.
\subsection{Stage Two: Untuned English-to-English Argument Comparison}
\draftspace
% TODO: original base with translation adapter disabled; same essential reasoning,
% assumptions, intermediate claims and dependencies rather than surface similarity.
\subsection{Evaluation Protocol and Direct-Judge Baseline}
\draftspace
% TODO: reserve identical cases for untuned direct English/Lean judge comparison.
% That matched-case baseline is next-stage work, not yet in the running queue.
% Report coverage; label provenance; mismatch P/R/F1, matrices and ROC only when
% correspondence labels justify them. Generation limits are not semantic NOs.

\section{Informalization Results and Targeted Case Studies}
% Budget: about 2 pages; primary emphasis, grading Experiments 40%.
\subsection{Held-Out Translation Quality and English Comparisons}
\draftspace
% TODO: discuss completed validation completion loss and example translations,
% judge decisions and truncation counts. Model judgments alone are not accuracy.
\begin{table*}[htbp]
\centering\small
\caption{Final informalization pipeline: completed observations and remaining evaluation.}
\begin{tabular}{lll}
\toprule
Evaluation & Evidence / metric & Status \\
\midrule
Herald translation & Validation NLL; held-out generations & 0.2948; 64/64 \\
Generation limits & Empty outputs / output-limit hits & 0 / 0 \\
English argument comparison & MATCH / MISMATCH / unscored & 49 / 15 / 0 \\
Matched-case direct judge & Same inputs and label policy & Next evaluation stage \\
OpenAI targeted cases & Reproduction of reported semantic discrepancies & Not yet run \\
Independent correspondence labels & Confusion matrix, F1, ROC; class coverage & Review required \\
\bottomrule
\end{tabular}
\end{table*}
\subsection{OpenAI Releases and Reported Navier--Stokes Discrepancies}
\draftspace
% TODO NEXT: locate and pin the exact released Lean/source-English cases cited
% by the critique and echoed by Xena. Record revisions, extraction and context.
% Test the final pipeline and direct baseline on identical cases. Include cases
% without reported discrepancies where feasible; lack of a report is not a gold
% positive label. Treat results as targeted case studies, not general accuracy.
% No OpenAI case-study job has been started or queued tonight.
\subsection{Original-Pair Errors Versus Informalizer and Judge Errors}
\draftspace
% TODO: inspect English source, Lean, generated English and judgment together;
% trace decisive steps back to Lean; avoid attributing all mismatches to source.
\subsection{Human Readability and Intermediate Representations}
\draftspace
% TODO: qualitative sanity-checking utility; controlled representations as future
% extensions. Embedding cosine is a proposed baseline, not a measured result.

\section{Limitations and Discussion}
% Budget: about 0.5 page.
\subsection{Sparse Public Labels, Shortcuts, and Evaluation Scope}
\draftspace
% TODO: label quality, small/repeated diagnostics, overlap and selected case studies.
\subsection{Model Scale, Workstation Resources, and Context Budgets}
\draftspace
% TODO: distinguish open-weight model capability from hardware and imposed token
% budgets. No frontier-model comparison was run; openness alone is not a cause.
% 32B inference worked; CPU/GPU backward probe stopped, not a tuned-32B result.

\section{Future Work: A Public Human-Verified Correspondence Dataset}
% Budget: about 0.4 page.
\subsection{Faithful and Unfaithful Proof Pairs with Step-Level Evidence}
\draftspace
% TODO: both labels within sources, same-theorem route changes, rationale,
% reviewer agreement/disagreement, whole-proof and step-level annotations.
\subsection{Pinned Formal Environments and Reusable Evaluation Splits}
\draftspace
% TODO: licenses, source provenance, Lean/mathlib versions, independently held-out
% theorem families; training judges and evaluating informalizers separately.

\section{Conclusion}
% Budget: about 0.2 page; grading weight 5%.
\draftspace
% TODO: Taylor's conclusions, based on completed evidence.

\bibliographystyle{ACM-Reference-Format}
% TODO: replace nocite with citations in your exposition; complete metadata.
\nocite{proofnet,herald,pistis,faithform,lora,qwen,navierstokes,xena,george}
\bibliography{references}
\end{document}
